% riccati.tex — §3–4 of the standalone paper
% The shadow metric, the Riccati algebraicity theorem, and the depth classification.

% ================================================================
% §3. The Shadow Metric and the Riccati Equation
% ================================================================

\section{The shadow metric and the Riccati equation}
\label{sec:shadow-metric}

\subsection{The shadow obstruction tower}

Let $\mathcal{A}$ be a modular Koszul chiral algebra with modular
cyclic deformation complex $\mathrm{Def}_{\mathrm{cyc}}^{\mathrm{mod}}(\mathcal{A})$
and universal Maurer--Cartan element
$\Theta_{\mathcal{A}} \in \mathrm{MC}(\mathrm{Def}_{\mathrm{cyc}}^{\mathrm{mod}}(\mathcal{A}))$.
The complex carries an arity filtration $F^{\bullet}$ with
$F^r = \bigoplus_{s \geq r} \mathrm{Def}_{\mathrm{cyc}}^{(s)}$,
and the \emph{shadow obstruction tower} is the sequence of
finite-order truncations
\[
 \Theta_{\mathcal{A}}^{\leq 2},\;
 \Theta_{\mathcal{A}}^{\leq 3},\;
 \Theta_{\mathcal{A}}^{\leq 4},\;
 \ldots
\]
where $\Theta_{\mathcal{A}}^{\leq r}$ solves the Maurer--Cartan equation
in the quotient $\mathfrak{g}_{\mathcal{A}}^{\mathrm{amb}} / F^{r+1}$.
The arity-$r$ projection $\mathrm{Sh}_r(\mathcal{A})$ is the
\emph{$r$-th shadow} of $\mathcal{A}$. On a one-dimensional
slice of the deformation complex, $\mathrm{Sh}_r$ is determined
by a single scalar $S_r$; the MC equation projected to arity $r$
gives the \emph{all-arity master equation}
\begin{equation}
\label{eq:master-equation}
 \nabla_H(\mathrm{Sh}_r) + \mathfrak{o}^{(r)} = 0,
\end{equation}
where $\nabla_H$ is the linearized operator and $\mathfrak{o}^{(r)}$
is the obstruction from lower arities.

\subsection{Primary lines and the shadow metric}

\begin{definition}[Primary line]
\label{def:primary-line}
A \emph{primary line} in $\mathrm{Def}_{\mathrm{cyc}}^{\mathrm{mod}}(\mathcal{A})$
is a one-dimensional subspace $L = \mathbf{k} \cdot \eta$
spanned by a single primary field $\eta$ of definite conformal
weight $h_\eta$. For a chiral algebra with generators
$\phi_1, \ldots, \phi_N$ of weights $h_1, \ldots, h_N$,
each generator spans a primary line; the full first-order
deformation space is $\bigoplus_i L_i$.
\end{definition}

Let $L$ be a primary line with coordinate $x$, and write
$\mathrm{Sh}_r|_L = S_r \, x^r$ for the shadow data restricted
to $L$. Set
\[
 \kappa := S_2, \qquad \alpha := S_3.
\]
The scalar $\kappa$ is the modular characteristic (curvature)
of $\mathcal{A}$ on~$L$; the scalar $\alpha$ is the cubic shadow.

\begin{definition}[Shadow metric]
\label{def:shadow-metric-paper}
Define the \emph{weighted initial data}
$a_0 := 2\kappa$, $a_1 := 3\alpha$, $a_2 := 4S_4$.
The \emph{shadow metric} (shadow quadratic form) on $L$ is the
polynomial
\begin{equation}
\label{eq:shadow-metric}
 Q_L(t)
 \;:=\;
 a_0^2 + 2\,a_0\,a_1\,t + (a_1^2 + 2\,a_0\,a_2)\,t^2
 \;=\;
 4\kappa^2 + 12\kappa\alpha\,t + (9\alpha^2 + 16\kappa S_4)\,t^2
 \;\in\; k(c)[t].
\end{equation}
The \emph{critical discriminant} of $L$ is
\begin{equation}
\label{eq:critical-discriminant}
 \Delta := 8\kappa S_4 = a_0 \, a_2.
\end{equation}
\end{definition}

\begin{remark}
The term ``metric'' is used by analogy with the role of $Q_L$
in governing the geometry of the shadow obstruction tower, not in the sense
of a Riemannian metric. The shadow metric is a quadratic
polynomial in the arity parameter $t$, determined by three OPE
invariants: $\kappa$, $\alpha$, and $S_4$.
\end{remark}

\subsection{The recursion on a primary line}

On a primary line $L$, the all-arity master equation
\eqref{eq:master-equation} reduces to a recursion among the scalars
$S_r$.

\begin{proposition}[Single-line recursion]
\label{prop:single-line-recursion}
On a primary line $L$ with propagator $P := 1/\kappa$,
the MC equation at arity $r \geq 4$ is
\begin{equation}
\label{eq:single-line-recursion}
 S_r
 \;=\;
 -\frac{P}{2r}
 \sum_{\substack{j + k = r + 2 \\ 3 \leq j \leq k}}
 c_{jk}\, j\, k\, S_j \, S_k,
 \qquad
 c_{jk} =
 \begin{cases}
 1 & \text{if } j < k, \\
 1/2 & \text{if } j = k.
 \end{cases}
\end{equation}
This recursion is well-defined: since $j \geq 3$ and $j + k = r + 2$,
every term on the right has $k = r + 2 - j \leq r - 1$, so $S_r$
is determined by $\{S_2, S_3, \ldots, S_{r-1}\}$.
\end{proposition}

\begin{proof}
The Lie bracket on the cyclic deformation complex, restricted to
$L$, pairs $\mathrm{Sh}_j = S_j x^j$ and
$\mathrm{Sh}_k = S_k x^k$ to produce a term of arity $j + k - 2$
with coefficient $jk\,P\,S_j S_k$. The eigenvalue identity
$\nabla_H(x^r) = 2r \, x^r$ separates the MC equation at
arity $r$ into the linear term $2r \, S_r$ on the left and the
obstruction sum on the right, giving the stated recursion.
The symmetry factor $c_{jk}$ accounts for the factor
$\tfrac{1}{2}$ in $\tfrac{1}{2}[\Theta, \Theta]$ when $j = k$.
\end{proof}

\subsection{The weighted shadow generating function}

\begin{definition}[Weighted shadow generating function]
\label{def:shadow-gf}
The \emph{weighted shadow generating function} on $L$ is
\begin{equation}
\label{eq:shadow-gf}
 H(t) \;:=\; \sum_{r \geq 2} r \, S_r \, t^r.
\end{equation}
Setting $a_n := (n+2)\,S_{n+2}$ for $n \geq 0$, this becomes
$H(t) = \sum_{n \geq 0} a_n \, t^{n+2}$, with initial values
$a_0 = 2\kappa$, $a_1 = 3\alpha$, $a_2 = 4S_4$.
\end{definition}

The weighted generating function $H$ packages the recursion
\eqref{eq:single-line-recursion} as a convolution identity.
Define $F(t) := H(t)/t^2 = \sum_{n \geq 0} a_n \, t^n$.
The Cauchy product $F^2 = \sum_{m \geq 0}
(\sum_{i+j=m} a_i a_j)\, t^m$ has its coefficients at
$m \geq 3$ controlled by the MC recursion. The next theorem
shows that the MC equation forces $F^2$ to be a polynomial
of degree $2$.


% ================================================================
% §4. Algebraicity and the Depth Classification
% ================================================================

\section{Algebraicity and the depth classification}
\label{sec:algebraicity}

\subsection{The Riccati algebraicity theorem}

\begin{theorem}[Riccati algebraicity]
\label{thm:riccati-algebraicity}
Let $L$ be a primary line with shadow data $S_r$, shadow metric
$Q_L$ \textup{(}Definition~\textup{\ref{def:shadow-metric-paper}}\textup{)},
and weighted shadow generating function $H$
\textup{(}Definition~\textup{\ref{def:shadow-gf}}\textup{)}.
Then the all-arity master equation on $L$ is equivalent to the
algebraic relation
\begin{equation}
\label{eq:riccati-algebraic-relation}
 H(t)^2 \;=\; t^4 \, Q_L(t).
\end{equation}
The shadow generating function is therefore algebraic of
degree $2$ over $k(\kappa, \alpha, S_4)(t)$:
\begin{equation}
\label{eq:shadow-closed-form}
 H(t) \;=\; t^2 \sqrt{Q_L(t)}.
\end{equation}
\end{theorem}

\begin{proof}
Set $F(t) := H(t)/t^2 = \sum_{n \geq 0} a_n \, t^n$ where
$a_n = (n+2)\, S_{n+2}$.
The claim \eqref{eq:riccati-algebraic-relation} is equivalent to
$F(t)^2 = Q_L(t)$. We verify this by computing the Cauchy
product $F^2 = \sum_{m \geq 0} \bigl(\sum_{i+j=m} a_i a_j\bigr)
\, t^m$ and showing it agrees with $Q_L$ at all orders.

\emph{Orders $0$, $1$, $2$ (initial data).}
At $m = 0$: $a_0^2 = (2\kappa)^2 = 4\kappa^2$.
At $m = 1$: $2\,a_0\,a_1 = 2 \cdot 2\kappa \cdot 3\alpha = 12\kappa\alpha$.
At $m = 2$: $a_1^2 + 2\,a_0\,a_2 = 9\alpha^2 + 2 \cdot 2\kappa \cdot 4S_4
= 9\alpha^2 + 16\kappa S_4$.
These are exactly the coefficients of $Q_L(t)$ in
\eqref{eq:shadow-metric}.

\emph{Orders $m \geq 3$ (the MC recursion forces vanishing).}
Since $Q_L$ has degree $2$ in $t$, the identity $F^2 = Q_L$ requires
\begin{equation}
\label{eq:convolution-vanishing}
 \sum_{i + j = m} a_i \, a_j \;=\; 0
 \qquad \text{for all } m \geq 3.
\end{equation}
We show that this is exactly the MC recursion
\eqref{eq:single-line-recursion} at arity $r = m + 2$.

The sum $\sum_{i+j=m} a_i a_j$ with
$a_i = (i+2)\,S_{i+2}$ and $a_j = (j+2)\,S_{j+2}$
equals $\sum_{p+q = m+4,\, p,q \geq 2} p\,q\,S_p\,S_q$,
where $p = i+2$, $q = j+2$. Setting $r = m+2$, this becomes
$\sum_{p+q = r+2,\, p,q \geq 2} p\,q\,S_p\,S_q$.

Isolating the terms with $p = 2$ or $q = 2$
(which contribute $2\kappa$ from $S_2 = \kappa$):
\begin{align*}
 \sum_{p+q=r+2,\, p,q \geq 2} pq\,S_p S_q
 &= 2 \cdot 2r \cdot \kappa\,S_r
 + \sum_{\substack{p+q=r+2 \\ p,q \geq 3}} pq\,S_p S_q \\
 &= 4r\kappa\,S_r
 + 2\!\!\sum_{\substack{j+k=r+2 \\ 3 \leq j \leq k}}\!\!
 c_{jk}\,jk\,S_j S_k,
\end{align*}
where the factor $2$ in front of the restricted sum accounts for
$S_p S_q + S_q S_p$ when $p \neq q$, and $c_{jk}$ corrects
for the diagonal. Substituting the recursion
\eqref{eq:single-line-recursion}:
\[
 2\!\!\sum_{\substack{j+k=r+2 \\ 3 \leq j \leq k}}\!\!
 c_{jk}\,jk\,S_j S_k
 \;=\;
 2 \cdot (-2r\kappa\,S_r)
 \;=\;
 -4r\kappa\,S_r.
\]
Hence $\sum_{p+q=r+2} pq\,S_p S_q = 4r\kappa S_r - 4r\kappa S_r = 0$,
which is \eqref{eq:convolution-vanishing}.

Conversely, the vanishing \eqref{eq:convolution-vanishing} at
order $m \geq 3$ implies the recursion
\eqref{eq:single-line-recursion} at $r = m+2$ by solving for
$S_r$ (using $\kappa \neq 0$).

The equivalence of \eqref{eq:riccati-algebraic-relation} with
the full MC recursion is therefore established at all arities.
Taking positive square roots (the branch with $F(0) = a_0 = 2\kappa > 0$)
gives $H(t) = t^2 F(t) = t^2 \sqrt{Q_L(t)}$.
\end{proof}

\begin{remark}[Riccati ODE reformulation]
\label{rem:riccati-ode}
Set $U(t) := \sum_{r \geq 3} S_r \, t^r$ (the nonlinear part,
removing $\mathrm{Sh}_2$). The algebraic relation
\eqref{eq:riccati-algebraic-relation} is equivalent to the
Riccati-type first-order ODE
\begin{equation}
\label{eq:riccati-ode}
 2t\,U'(t) + \frac{P}{2}\bigl(U'(t)\bigr)^2
 = R_3\,t^3 + R_4\,t^4,
\end{equation}
where $R_3 = 6\alpha$ and
$R_4 = (9\alpha^2 + 2\Delta)/(2\kappa)$ encode the cubic
and quartic inputs. The left-hand side is quadratic in $U'$:
this is the precise sense in which the shadow obstruction tower is algebraic
of degree $2$.
\end{remark}

\subsection{Gaussian decomposition}

\begin{corollary}[Gaussian decomposition]
\label{cor:gaussian-decomposition}
The shadow metric decomposes as
\begin{equation}
\label{eq:gaussian-decomposition}
 Q_L(t) \;=\; (2\kappa + 3\alpha\,t)^2 + 2\Delta\,t^2,
\end{equation}
where $\Delta = 8\kappa S_4$ is the critical discriminant.
The first term $(2\kappa + 3\alpha t)^2$ is the \emph{Gaussian
envelope}: the shadow metric of the truncated tower with
$S_r = 0$ for $r \geq 4$. The second term $2\Delta\,t^2$
is the \emph{interaction correction}, measuring the quartic
OPE contribution not captured by the cubic data alone.
\end{corollary}

\begin{proof}
$(2\kappa + 3\alpha t)^2 = 4\kappa^2 + 12\kappa\alpha\,t
+ 9\alpha^2\,t^2$. Adding $2\Delta\,t^2 = 16\kappa S_4\,t^2$
recovers the coefficient $9\alpha^2 + 16\kappa S_4$ of $t^2$
in $Q_L(t)$.
\end{proof}

\subsection{Intrinsic quartic principle}

\begin{corollary}[Intrinsic quartic principle]
\label{cor:intrinsic-quartic}
On a primary line $L$, the shadow obstruction tower has intrinsic OPE data
at arities $2$, $3$, $4$ only. All shadow coefficients $S_r$
for $r \geq 5$ are determined by $(\kappa, \alpha, S_4)$
alone, via the closed form
\begin{equation}
\label{eq:shadow-from-closed-form}
 S_r \;=\; \frac{1}{r}\,[t^{r-2}]\,\sqrt{Q_L(t)}
 \qquad (r \geq 2).
\end{equation}
\end{corollary}

\begin{proof}
$Q_L(t) = 4\kappa^2 + 12\kappa\alpha\,t + (9\alpha^2 + 16\kappa S_4)\,t^2$
depends on exactly $(\kappa, \alpha, S_4)$.
Since $H(t) = t^2\sqrt{Q_L(t)} = \sum_{r \geq 2} rS_r\,t^r$
by Theorem~\ref{thm:riccati-algebraicity}, comparing
coefficients gives $rS_r = [t^{r-2}]\sqrt{Q_L}$.
\end{proof}

\subsection{The shadow connection}

\begin{proposition}[Shadow connection]
\label{prop:shadow-connection}
The shadow metric defines a logarithmic connection on
$\mathcal{O}_L$ over $L \setminus \{Q_L = 0\}$:
\begin{equation}
\label{eq:shadow-connection}
 \nabla^{\mathrm{sh}} \;=\; d - \omega,
 \qquad
 \omega \;:=\; \frac{Q_L'(t)}{2\,Q_L(t)}\,dt
 \;=\; d\bigl(\log\sqrt{Q_L}\bigr).
\end{equation}
\begin{enumerate}[label=\textup{(\roman*)}]
\item The connection is flat: $(\nabla^{\mathrm{sh}})^2 = 0$.
\item At each simple zero $t_0$ of $Q_L$, the residue
 of $\omega$ is $\tfrac{1}{2}$.
\item The monodromy is $\exp(2\pi i \cdot \tfrac{1}{2}) = -1$:
 the standard monodromy of a square root around a simple
 branch point.
\item The flat section with $\Phi(0) = 1$ is
 $\Phi(t) = \sqrt{Q_L(t)/Q_L(0)} = \sqrt{Q_L(t)}/(2\kappa)$,
 and the shadow generating function is
 $H(t) = 2\kappa\,t^2\,\Phi(t)$: the parallel transport of
 the curvature along the arity parameter, weighted by $t^2$.
\end{enumerate}
\end{proposition}

\begin{proof}
Flatness is automatic: $\omega$ is a closed $1$-form on a
$1$-dimensional base. At a simple zero $t_0$ of $Q_L$:
$Q_L(t) \approx Q_L'(t_0)(t - t_0)$, so
$\omega \approx \frac{1}{2(t - t_0)}\,dt$,
giving residue $\tfrac{1}{2}$. The monodromy exponentiation
is immediate. For (iv): $\nabla^{\mathrm{sh}}\Phi = 0$
since $d\Phi = \omega\,\Phi$ by construction, and the
evaluation at $t = 0$ gives
$\Phi(0) = \sqrt{Q_L(0)}/(2\kappa) = 2\kappa/(2\kappa) = 1$.
\end{proof}

\begin{remark}
The shadow generating function $H(t) = 2\kappa\,t^2\,\Phi(t)$
is \emph{not} a flat section of $\nabla^{\mathrm{sh}}$: the
factor $t^2$ contributes $\nabla^{\mathrm{sh}}(H) \neq 0$.
The flat section is $\Phi(t) = \sqrt{Q_L(t)}/(2\kappa)$; the
shadow obstruction tower is the Taylor expansion of $t^2 \cdot \Phi(t)$,
reflecting the arity offset (the tower starts at arity $2$,
not arity $0$).
\end{remark}

\subsection{The single-line dichotomy}

\begin{theorem}[Single-line dichotomy]
\label{thm:single-line-dichotomy}
Let $L$, $S_r$, $\alpha$, $\Delta$, $Q_L$ be as above,
and define the shadow depth $r_{\max}|_L := \max\{r : S_r \neq 0\}$
\textup{(}with $r_{\max} = \infty$ if $S_r \neq 0$ for
infinitely many $r$\textup{)}.
Then $r_{\max}|_L \in \{2, 3, \infty\}$, classified by the
shadow metric:
\begin{enumerate}[label=\textup{(\roman*)}]
\item \textbf{Class $\mathbf{G}$}
 \textup{(}$\Delta = 0$, $\alpha = 0$\textup{)}.\enspace
 $Q_L = (2\kappa)^2$ is constant.
 $H(t) = 2\kappa\,t^2$. Only $S_2$ survives: $r_{\max} = 2$.

\item \textbf{Class $\mathbf{L}$}
 \textup{(}$\Delta = 0$, $\alpha \neq 0$\textup{)}.\enspace
 $Q_L = (2\kappa + 3\alpha\,t)^2$ is a perfect square.
 $H(t) = t^2(2\kappa + 3\alpha\,t)$. $S_2$ and $S_3$
 survive: $r_{\max} = 3$.

\item \textbf{Class $\mathbf{M}$}
 \textup{(}$\Delta \neq 0$\textup{)}.\enspace
 $Q_L$ is irreducible over $k(c)$; $\sqrt{Q_L}$ is
 irrational; $S_r \neq 0$ for all $r \geq 4$ generically:
 $r_{\max} = \infty$.
\end{enumerate}
For $r \geq 4$,
\begin{equation}
\label{eq:universal-factorization}
 S_r = \Delta \cdot R_r,
 \qquad R_r \in \mathbb{Q}(\alpha, \kappa, S_4).
\end{equation}
\end{theorem}

\begin{proof}
The classification follows from Theorem~\ref{thm:riccati-algebraicity}
and the Gaussian decomposition
(Corollary~\ref{cor:gaussian-decomposition}).

\emph{Case $\Delta = 0$.}
Then $Q_L = (2\kappa + 3\alpha t)^2$ is a perfect square, so
$\sqrt{Q_L} = 2\kappa + 3\alpha t$ is polynomial and
$H(t) = t^2(2\kappa + 3\alpha t)$ terminates.
If $\alpha = 0$: $H(t) = 2\kappa t^2$, giving $r_{\max} = 2$.
If $\alpha \neq 0$: $H(t) = 2\kappa t^2 + 3\alpha t^3$,
giving $r_{\max} = 3$.

\emph{Case $\Delta \neq 0$.}
The discriminant of $Q_L$ as a polynomial in $t$ is
\[
 (12\kappa\alpha)^2 - 4 \cdot 4\kappa^2 \cdot
 (9\alpha^2 + 16\kappa S_4) = -32\kappa^2\Delta \neq 0,
\]
so $Q_L$ is irreducible over $k(c)$. The Taylor expansion
of $\sqrt{Q_L}$ around $t = 0$ has nonzero coefficients at
all orders generically, hence $r_{\max} = \infty$.

\emph{Universal factorization.}
The binomial expansion of $\sqrt{Q_L}$ around the Gaussian
envelope $(2\kappa + 3\alpha t)$ is
\[
 \sqrt{Q_L} = (2\kappa + 3\alpha t)\sqrt{1 + \frac{2\Delta\,t^2}
 {(2\kappa + 3\alpha t)^2}}.
\]
Expanding $\sqrt{1 + u}$ as $1 + \tfrac{1}{2}u - \tfrac{1}{8}u^2
+ \cdots$ with $u = 2\Delta t^2/(2\kappa + 3\alpha t)^2$,
every coefficient of $t^n$ for $n \geq 2$ in
$\sqrt{Q_L} - (2\kappa + 3\alpha t)$ carries a factor of
$\Delta$. Since $rS_r = [t^{r-2}]\sqrt{Q_L}$ for $r \geq 2$,
and the Gaussian part $2\kappa + 3\alpha t$ contributes only
at $r \leq 3$, the factorization $S_r = \Delta \cdot R_r$
for $r \geq 4$ follows.
\end{proof}

\begin{remark}[The contact class]
\label{rem:contact-class}
The partition $\mathbf{G}/\mathbf{L}/\mathbf{C}/\mathbf{M}$
of modular Koszul algebras includes a fourth class $\mathbf{C}$
(contact, $r_{\max} = 4$), realized by the $\beta\gamma$ system.
This class escapes the single-line dichotomy by \emph{stratum
separation}: the quartic contact invariant lives on a charged
stratum of $\mathrm{Def}_{\mathrm{cyc}}^{\mathrm{mod}}$ where
$\kappa = 0$, and the shadow metric requires $\kappa \neq 0$.
On the charged stratum, the self-bracket maps to a zero-dimensional
target by rank-one rigidity, so the quartic pump does not activate
and the tower terminates at $r = 4$. The algebraic shadow depth
$d_{\mathrm{alg}} \in \{0, 1, 2, \infty\}$ exhausts the four
classes: $\mathbf{G}$ ($d_{\mathrm{alg}} = 0$),
$\mathbf{L}$ ($d_{\mathrm{alg}} = 1$),
$\mathbf{C}$ ($d_{\mathrm{alg}} = 2$),
$\mathbf{M}$ ($d_{\mathrm{alg}} = \infty$).
\end{remark}

\subsection{The shadow growth rate}

For class $\mathbf{M}$ algebras, the shadow obstruction tower is infinite
but governed by a single analytic invariant: the growth rate of
the coefficients $S_r$.

\begin{definition}[Shadow growth rate]
\label{def:shadow-growth-rate}
Let $L$ be a primary line of class $\mathbf{M}$ ($\Delta \neq 0$),
and write $Q_L(t) = q_0 + q_1 t + q_2 t^2$ for the shadow metric.
The \emph{shadow growth rate} of $L$ is
\begin{equation}
\label{eq:shadow-growth-rate}
 \rho_L \;:=\; \sqrt{\frac{q_2}{q_0}}
 \;=\;
 \frac{\sqrt{9\alpha^2 + 2\Delta}}{2|\kappa|}.
\end{equation}
By Vieta's theorem, the zeros $t_0$, $\bar{t}_0$ of $Q_L$
satisfy $|t_0|^2 = q_0/q_2$, so
$\rho_L = 1/|t_0|$: the reciprocal of the branch-point modulus.
For classes $\mathbf{G}$ and $\mathbf{L}$, set $\rho := 0$
by convention.
\end{definition}

\begin{theorem}[Shadow asymptotics]
\label{thm:shadow-asymptotics}
Let $L$ be a primary line of class $\mathbf{M}$ with shadow
growth rate $\rho_L$.
\begin{enumerate}[label=\textup{(\roman*)}]
\item \emph{Asymptotic formula.}
 \begin{equation}
 \label{eq:shadow-asymptotic-formula}
 S_r \;\sim\; A(\mathcal{A})\,\rho_L^{\,r}\,r^{-5/2}\,
 \cos(r\theta + \varphi)
 \qquad (r \to \infty),
 \end{equation}
 where $\theta := -\arg(t_0)$ is the oscillation phase,
 $A(\mathcal{A}) > 0$, and $\varphi$ is computable from the
 branch-point residue of $\sqrt{Q_L}$. The exponent $-5/2$
 arises from $S_r = [t^{r-2}]\sqrt{Q_L}/r$ combined with the
 $n^{-3/2}$ transfer law for square-root singularities.

\item \emph{Unit circle criterion.}
 $\rho_L = 1$ if and only if $q_0 = q_2$, i.e.,
 \begin{equation}
 \label{eq:unit-circle-criterion}
 (2\kappa - 3\alpha)(2\kappa + 3\alpha) = 2\Delta.
 \end{equation}
 Equivalently, the branch points of $Q_L$ lie on the unit
 circle $|t| = 1$.

\item \emph{Convergence partition.}
 $\rho_L < 1$: the shadow obstruction tower converges
 absolutely.
 $\rho_L > 1$: the shadow obstruction tower diverges, but the algebraic
 function $\sqrt{Q_L(t)}$ extends analytically past the
 convergence radius.
 $\rho_L = 1$: marginal case (the critical surface).
\end{enumerate}
\end{theorem}

\begin{proof}
The function $\sqrt{Q_L(t)}$ has algebraic branch-cut
singularities of order $\tfrac{1}{2}$ at each simple zero
$t_0$ of $Q_L$. Since $Q_L$ is quadratic with nonzero
discriminant ($\Delta \neq 0$ implies
$\mathrm{disc}(Q_L) = -32\kappa^2\Delta \neq 0$),
the two zeros $t_0, \bar{t}_0$ are distinct with
$|t_0| = |{\bar t}_0| = \sqrt{q_0/q_2} = 1/\rho_L$.

The transfer theorem for algebraic singularities
(Flajolet--Sedgewick, Theorem~VI.1) gives
$[t^n]\sqrt{Q_L} \sim C_0\,t_0^{-n}\,n^{-3/2}$
for each dominant singularity $t_0$ at distance
$R = 1/\rho_L$ from the origin. Both branch points
contribute (equal modulus), producing the cosine
interference term from
$2\operatorname{Re}(C_0\,t_0^{-(r-2)})$.
Since $S_r = [t^{r-2}]\sqrt{Q_L}/r$
(Corollary~\ref{cor:intrinsic-quartic}), the extra
factor $1/r$ gives the stated $r^{-5/2}$ exponent.

For (ii): $\rho_L = 1$ iff $q_2/q_0 = 1$. Writing
$q_0 = 4\kappa^2$ and $q_2 = 9\alpha^2 + 2\Delta$, the
condition $q_0 = q_2$ factors as
$4\kappa^2 - 9\alpha^2 = 2\Delta$.
Part (iii) is immediate from the definition of the
convergence radius $R = 1/\rho_L$.
\end{proof}

\subsection{The depth decomposition}

\begin{theorem}[Depth decomposition]
\label{thm:depth-decomposition}
The shadow depth of a modular Koszul algebra $\mathcal{A}$
decomposes as
\begin{equation}
\label{eq:depth-decomposition}
 d(\mathcal{A}) \;=\; 1 + d_{\mathrm{arith}}(\mathcal{A})
 + d_{\mathrm{alg}}(\mathcal{A}),
\end{equation}
where:
\begin{enumerate}[label=\textup{(\roman*)}]
\item $d_{\mathrm{arith}}$ counts the independent holomorphic
 Hecke eigenforms in the spectral decomposition of the completed
 shadow partition function $\widehat{Z}^c_{\mathcal{A}}$
 on $\mathcal{M}_{1,1}$;

\item $d_{\mathrm{alg}} \in \{0, 1, 2, \infty\}$ is the
 \emph{algebraic depth}, determined by the single-line dichotomy
 \textup{(}Theorem~\textup{\ref{thm:single-line-dichotomy}}\textup{)}
 and stratum separation
 \textup{(}Remark~\textup{\ref{rem:contact-class}}\textup{)}:
 \[
 d_{\mathrm{alg}} =
 \begin{cases}
 0 & \text{class } \mathbf{G}, \\
 1 & \text{class } \mathbf{L}, \\
 2 & \text{class } \mathbf{C}, \\
 \infty & \text{class } \mathbf{M};
 \end{cases}
 \]

\item depths $d \geq 5$ are purely arithmetic: they arise from
 $d_{\mathrm{arith}}$, not from the $A_\infty$ structure.
\end{enumerate}
\end{theorem}

\begin{proof}
The algebraic depth $d_{\mathrm{alg}}$ is determined by the
structure of the cyclic deformation complex and the MC recursion:
the single-line dichotomy
(Theorem~\ref{thm:single-line-dichotomy}) gives
$d_{\mathrm{alg}} \in \{0, 1, \infty\}$ on each primary line,
and stratum separation (Remark~\ref{rem:contact-class}) provides
$d_{\mathrm{alg}} = 2$ for the contact class.
Since $d_{\mathrm{alg}} \leq 2$ for finite cases
($\mathbf{G}/\mathbf{L}/\mathbf{C}$) and the base term
is $1$ (the curvature $\kappa$ at arity $2$), the maximum
finite depth without arithmetic contributions is
$1 + 0 + 2 = 3$; hence all depths $d \geq 5$ require
$d_{\mathrm{arith}} \geq 2$.

The arithmetic depth counts independent holomorphic Hecke
eigenforms by decomposing $\widehat{Z}^c_{\mathcal{A}}$
via the Roelcke--Selberg spectral decomposition on
$\mathcal{M}_{1,1}$. Each holomorphic eigenform of weight $k$
contributes a critical line at $\operatorname{Re}(s) = k/2$;
Maass cusp-form projections contribute no new critical lines.
\end{proof}

\begin{example}[Standard families]
\label{ex:standard-families}
The shadow metric data for the standard families:
\begin{center}
\small
\renewcommand{\arraystretch}{1.3}
\begin{tabular}{lcccccl}
\toprule
Family & $\kappa$ & $\alpha$ & $S_4$ & $\Delta$ & Class
 & $\rho$ \\
\midrule
Heisenberg $\mathcal{H}_k$ & $k$ & $0$ & $0$ & $0$
 & $\mathbf{G}$ & $0$ \\
Affine $\widehat{\mathfrak{g}}_k$
 & $\frac{d(k+h^\vee)}{2h^\vee}$ & $\neq 0$ & $0$
 & $0$ & $\mathbf{L}$ & $0$ \\
$\beta\gamma$ (weight $\lambda$) & $6\lambda^2{-}6\lambda{+}1$ & $0$ & $\neq 0$ & $0^*$
 & $\mathbf{C}$ & $*$ \\
Virasoro $\mathrm{Vir}_c$ & $c/2$ & $2$
 & $\frac{10}{c(5c+22)}$ & $\frac{40}{5c+22}$
 & $\mathbf{M}$ & $\frac{1}{c}\sqrt{\frac{180c+872}{5c+22}}$ \\
\bottomrule
\end{tabular}
\end{center}
\vspace{2pt}
\noindent
${}^*$\,Stratum separation: $\Delta = 0$ on the primary line
(where $S_4 = 0$); the quartic contact invariant
$Q^{\mathrm{contact}} \neq 0$ lives on a charged stratum
where $\kappa = 0$, and the single-line dichotomy does not
apply.

For the Virasoro algebra, the shadow metric evaluates to
\[
 Q_{\mathrm{Vir}}(t) = c^2 + 12c\,t
 + \frac{180c + 872}{5c + 22}\,t^2,
\]
with Gaussian decomposition
$Q_{\mathrm{Vir}} = (c + 6t)^2 + \frac{80}{5c+22}\,t^2$.
The critical central charge $c^*$ where $\rho = 1$ is the
unique positive real root of
\[
 5c^3 + 22c^2 - 180c - 872 = 0,
\]
approximately $c^* \approx 6.12$. For $c > c^*$ the shadow
tower converges; for $c < c^*$ it diverges. At the Koszul
self-dual point $c = 13$: $\rho(\mathrm{Vir}_{13}) \approx 0.467$.
\end{example}

\begin{remark}[Spectral curve]
The algebraic relation $H^2 = t^4 Q_L(t)$ defines, for each
primary line $L$, a spectral curve
$\Sigma_L := \{H^2 = t^4 Q_L(t)\} \subset \mathbb{A}^2_{t,H}$.
Since $Q_L$ is quadratic in $t$, $\Sigma_L$ has arithmetic
genus $0$; the projection $(t,H) \mapsto t$ is a double cover
ramified at the zeros of $Q_L$. The degeneration loci in
parameter space are controlled by the critical discriminant:
$\Sigma_L$ acquires a node precisely when $\Delta = 0$.
The shadow growth rate $\rho_L$ is the reciprocal of the
distance from $t = 0$ to the nearest ramification point.
\end{remark}
